\label{sec:introduction}

Gerrymandering's most durable electoral consequence is the elimination of
competitive congressional races.  When mapmakers can see partisan data at
the census-block level and adjust boundaries in real time to achieve desired
seat outcomes, safe districts are the almost inevitable result.  The enacted
maps from the 2020 redistricting cycle illustrate the point starkly: of the
435 congressional districts drawn for the 117th Congress, fewer than 15\%
were decided by margins under 10 percentage points, and the median margin of
victory exceeded 30 points~\citep{jacobson2012politics}.

Against this backdrop, one argument frequently advanced for algorithmic
redistricting is that geometrically neutral maps will---as a byproduct of
their neutrality---produce more competitive elections.  If the algorithm
cannot see partisan data, it cannot crack and pack partisans into safe seats.
Competitive districts should emerge naturally from the geometry of where
people live.

This paper tests that argument rigorously.  We apply 2020 presidential vote
returns to algorithmically generated congressional districts for all 50 states
and compare the distribution of competitive races to that obtained under
enacted plans.  We do not claim that the algorithm \emph{targets}
competitiveness---it does not---but we do claim that competitiveness is a
measurable byproduct of geometric neutrality.

The democratic-responsiveness argument for electoral competition has both
empirical and normative dimensions.  Empirically, competitive districts
produce higher turnout, stronger constituency service, and greater
policy responsiveness~\citep{jacobson2015it}.  They force incumbents to build
genuine majority coalitions rather than mobilize a safe partisan base.  They
also reduce the ``demand side'' of negative partisanship: voters in
uncompetitive districts disengage from electoral politics partly because
outcomes are foregone conclusions~\citep{abramowitz2006incumbency}.

Normatively, competitive districts are often identified with the ideal of
electoral accountability: voters should be able to throw incumbents out if
they fail to represent constituent interests.  In a legislature where almost
every seat is safe, accountability at the district level is impaired even if
national-level partisan tides occasionally flip chamber control.

A counterargument, associated with \citet{brunell2008redistricting}, holds that
safe districts are normatively preferable because they maximize the share of
constituents represented by their preferred party.  We engage this view in
our discussion of partisan asymmetry.  Even granting Brunell's normative
premises, we show that competitive algorithmic districts have a near-symmetric
partisan split---the claim of systematic Democratic or Republican advantage
under algorithmic maps does not survive empirical scrutiny.

The remainder of the paper proceeds as follows.  Section~2 defines
competition operationally.  Section~3 describes the data and methodology.
Section~4 presents the main results.  Sections~5--7 analyze partisan
asymmetry, longitudinal stability, and state heterogeneity.  Section~8
concludes.
